% !TeX root = ../main.tex
\chapter{Detailed Software Flowchart}
\label{app:software-flowchart}
This appendix provides the detailed software flow requested for the Week 2
preliminary-design package. It shows the firmware acquisition/alarm loop and the
PSC-to-EPICS data path used by the implemented software.

The software for the Multi-Rail Isolated Power Supply Monitoring System is
maintained with the hardware design files in the project Git repository:

\begin{center}
    \url{https://github.com/capotosto/PS_Monitor}
\end{center}

The repository provides version control for the Arduino GIGA firmware,
hardware-design files, and EPICS IOC software used by the monitoring system.
The embedded firmware is maintained under the \texttt{firmware} portion of the
repository, while the EPICS IOC and associated database/configuration files are
maintained under \texttt{epics-ioc}. The flowcharts in this section describe
the major software operations implemented by these portions of the project.

The Arduino GIGA firmware performs system initialization, acquires voltage and
current measurements from the six LTC2945 monitor channels, evaluates alarm
conditions, updates the local display, and publishes data through the Ethernet interface. The primary firmware execution
sequence is:

\begin{enumerate}
    \item Initialize the GIGA Display, persistent settings, Ethernet interface,
    I\textsuperscript{2}C interface, and channel configuration.

    \item Load stored network parameters and voltage/current alarm limits from
    persistent storage.

    \item Initialize communication with each of the six LTC2945 monitoring
    channels.

    \item Enter the continuous acquisition and communication loop.

    \item Poll each LTC2945 channel for voltage and current measurements.

    \item Verify I\textsuperscript{2}C communication status for each channel.

    \item If communication fails, mark the associated channel with an
    I\textsuperscript{2}C error and prevent invalid measurement data from being
    used.

    \item Convert valid raw measurement data to engineering units and apply the
    appropriate channel polarity and calibration/scaling information.

    \item Compare each voltage and current measurement with its configured
    upper and lower alarm limits.

    \item Update the GIGA Display when measurement values, alarm states, or
    communication states change.

    \item Publish the current measurements and status information through the
    Ethernet PSC protocol interface.

    \item Service incoming configuration requests and store modified settings
    when required.

    \item Repeat the acquisition and communication sequence continuously.
\end{enumerate}


% =====================================================================
\section{Arduino GIGA Firmware Flowchart}
\label{sec:giga-firmware-flowchart}

Figure~\ref{fig:firmware-flowchart} shows the top-level execution flow of the
Arduino GIGA firmware. Following initialization, the controller continuously
acquires data from each monitoring channel, verifies communication, converts
valid measurements to engineering units, evaluates alarm limits, updates the
local display, and services the Ethernet/PSC interface.

\begin{figure}[H]
    \centering

    \begin{adjustbox}{
            max totalsize={0.94\textwidth}{0.88\textheight},
            center
        }

        \begin{tikzpicture}[
            >=Stealth,
            font=\small,
            node distance=0.55cm and 1.6cm,
            process/.style={
                draw,
                rectangle,
                rounded corners=2pt,
                align=center,
                minimum width=5.0cm,
                minimum height=0.80cm,
                text width=4.7cm,
                inner sep=4pt,
                line width=0.8pt
            },
            decision/.style={
                draw,
                diamond,
                aspect=2.4,
                align=center,
                text width=2.8cm,
                inner sep=1pt,
                line width=0.8pt
            },
            terminator/.style={
                draw,
                rectangle,
                rounded corners=8pt,
                align=center,
                minimum width=3.2cm,
                minimum height=0.75cm,
                line width=0.8pt
            },
            arrow/.style={
                ->,
                line width=0.9pt
            },
            flowlabel/.style={
                font=\scriptsize,
                fill=white,
                inner sep=1.5pt
            }
            ]

            % ============================================================
            % INITIALIZATION
            % ============================================================

            \node[terminator] (start)
            {Power-Up / Reset};

            \node[process, below=of start] (init)
            {Initialize GIGA Hardware\\
                Display, I\textsuperscript{2}C, Ethernet, and GPIO};

            \node[process, below=of init] (settings)
            {Load Persistent Configuration\\
                Network Settings and Alarm Limits};

            \node[process, below=of settings] (configure)
            {Configure Six LTC2945 Monitor Channels\\
                and Initialize PSC Communication};

            % ============================================================
            % ACQUISITION LOOP
            % ============================================================

            \node[process, below=0.9cm of configure] (select)
            {Select Next Monitor Channel};

            \node[process, below=of select] (read)
            {Request Voltage and Current\\
                Measurements from LTC2945};

            \node[decision, below=0.70cm of read] (comm)
            {I\textsuperscript{2}C Communication\\Successful?};

            % Error branch
            \node[process, right=2.1cm of comm] (error)
            {Set Channel\\
                I\textsuperscript{2}C Error Status};

            % Valid measurement path
            \node[process, below=0.75cm of comm] (convert)
            {Convert Raw ADC Data to Engineering Units\\
                and Apply Scaling / Rail Polarity};

            \node[process, below=of convert] (alarms)
            {Compare Voltage and Current\\
                Against Configured Alarm Limits};

            \node[decision, below=0.70cm of alarms] (last)
            {Last Channel?};

            % ============================================================
            % SYSTEM UPDATE
            % ============================================================

            \node[process, below=0.75cm of last] (display)
            {Update Changed Display Fields\\
                Measurements, Alarms, and Status};

            \node[process, below=of display] (psc)
            {Publish Measurements and Status\\
                through PSC / Ethernet};

            \node[decision, below=0.70cm of psc] (configreq)
            {Configuration\\Request Received?};

            \node[process, right=2.1cm of configreq] (save)
            {Process Request and Save\\
                Persistent Settings};

            % ============================================================
            % PRIMARY FLOW
            % ============================================================

            \draw[arrow] (start) -- (init);
            \draw[arrow] (init) -- (settings);
            \draw[arrow] (settings) -- (configure);
            \draw[arrow] (configure) -- (select);
            \draw[arrow] (select) -- (read);
            \draw[arrow] (read) -- (comm);

            % Communication successful
            \draw[arrow]
            (comm.south)
            --
            node[flowlabel, right] {Yes}
            (convert.north);

            \draw[arrow] (convert) -- (alarms);
            \draw[arrow] (alarms) -- (last);

            % Communication failure
            \draw[arrow]
            (comm.east)
            --
            node[flowlabel, above] {No}
            (error.west);

            % Error skips conversion/alarm processing
            \draw[arrow]
            (error.south)
            |- (last.east);

            % More channels remain
            \draw[arrow]
            (last.west)
            -- ++(-1.3,0)
            |- node[flowlabel, near start, above] {No}
            (select.west);

            % All six channels complete
            \draw[arrow]
            (last.south)
            --
            node[flowlabel, right] {Yes}
            (display.north);

            \draw[arrow] (display) -- (psc);
            \draw[arrow] (psc) -- (configreq);

            % Configuration request
            \draw[arrow]
            (configreq.east)
            --
            node[flowlabel, above] {Yes}
            (save.west);

            % No configuration request:
            % loop back without dropping below the figure
            \coordinate (loopbottom)
            at ($(configreq.south)+(0,-0.45)$);

            \draw[arrow]
            (configreq.south)
            --
            node[flowlabel, right] {No}
            (loopbottom)
            -| ($(select.west)+(-1.3,0)$)
            |- (select.west);

            % After saving configuration, re-enter acquisition loop
            \draw[arrow]
            (save.south)
            |- ($(configreq.south)+(0,-0.45)$);

        \end{tikzpicture}

    \end{adjustbox}

    \caption{
        Top-level Arduino GIGA firmware flowchart. Following initialization,
        the controller continuously acquires measurements from the six LTC2945
        monitoring channels, validates I\textsuperscript{2}C communication,
        converts valid data to engineering units, evaluates configured alarm
        limits, updates the local display, publishes measurements through the
        PSC/Ethernet interface, and services configuration requests.
    }
    \label{fig:firmware-flowchart}

\end{figure}


\section{PSC and EPICS Communication Flow}
\label{sec:epics-software-flowchart}
% =====================================================================

The Arduino GIGA does not directly implement EPICS Channel Access. Instead,
measurement and status information is transmitted through the Ethernet
interface using the PSC protocol. The EPICS IOC uses PSCDriver to communicate
with the GIGA, maps the received values into EPICS process variables, and makes
those process variables available to the CS-Studio Phoebus operator interface.

This separation allows the embedded controller to remain responsible for
hardware acquisition and local alarm evaluation while the IOC provides the
interface between the embedded hardware and the facility control system.
Figure~\ref{fig:epics-data-flow} summarizes this software data path.

\begin{figure}[H]
    \centering

    \resizebox{0.95\textwidth}{!}{%
        \begin{tikzpicture}[
            >=Stealth,
            node distance=1.25cm,
            process/.style={
                draw,
                rectangle,
                rounded corners=2pt,
                align=center,
                minimum width=3.2cm,
                minimum height=1.5cm,
                text width=3.0cm,
                line width=0.8pt
            },
            arrow/.style={
                ->,
                line width=1.0pt
            },
            siglabel/.style={
                font=\scriptsize,
                fill=white,
                inner sep=2pt,
                rotate=90
            }
            ]

            \node[process] (measure)
            {Arduino GIGA\\[2pt]
                \footnotesize
                Acquire and Process\\
                Measurements};

            \node[process, right=1.5cm of measure] (psc)
            {PSC Protocol\\[2pt]
                \footnotesize
                Package Voltage, Current,\\
                Alarm, and Status Data};

            \node[process, right=1.5cm of psc] (ethernet)
            {Ethernet Interface\\[2pt]
                \footnotesize
                Network Transport};

            \node[process, right=1.5cm of ethernet] (ioc)
            {EPICS IOC / PSCDriver\\[2pt]
                \footnotesize
                Decode Data and Update\\
                EPICS Process Variables};

            \node[process, right=1.5cm of ioc] (phoebus)
            {CS-Studio Phoebus\\[2pt]
                \footnotesize
                Display Measurements,\\
                Alarms, and Status};

            \draw[arrow]
            (measure.east) --
            node[siglabel, midway] {Processed Data}
            (psc.west);

            \draw[arrow]
            (psc.east) --
            node[siglabel, midway] {PSC Messages}
            (ethernet.west);

            \draw[arrow]
            (ethernet.east) --
            node[siglabel, midway] {Ethernet}
            (ioc.west);

            \draw[arrow]
            (ioc.east) --
            node[siglabel, midway] {EPICS Channel Access}
            (phoebus.west);

        \end{tikzpicture}%
    }

    \caption{
        Software communication path from the Arduino GIGA firmware through
        PSCDriver and the EPICS IOC to the CS-Studio Phoebus operator interface.
    }
    \label{fig:epics-data-flow}
\end{figure}


% =====================================================================
\section{Software Configuration Management}
\label{sec:software-configuration}
% =====================================================================

Firmware, IOC configuration files, database definitions, and operator-interface
files are maintained under Git version control so that hardware and software
revisions can be associated with a defined project baseline. The current
software implementation and revision history are available in the
\texttt{PS\_Monitor} repository at:

\begin{center}
    \url{https://github.com/capotosto/PS_Monitor}
\end{center}

Changes made during implementation and testing are committed to the repository,
providing traceability between the preliminary design presented in this report
and the final implemented system. The repository therefore serves as the
controlled source for the firmware and control-system software associated with
the project.

